% BLOCK12_STAGE2_MANAGED
% Lecture 12: move semantics, value categories, forwarding and explicit casts
\section{Лекция 12. Move-семантика и явные приведения типов}
\label{sec:lecture-12}

В лекции~\ref{sec:lecture-11} возникла странная на первый взгляд операция:
\texttt{unique\_ptr} нельзя копировать, но его можно передать другому owner через
\texttt{std::move}. Теперь разберём, что именно означает эта запись и почему
C++ различает выражения не только по типу, но и по \term{value category}.

Вторая половина лекции посвящена похожей design-идее: если programmer намеренно
меняет способ интерпретации значения или типа, это намерение лучше показать
явно. Поэтому современный C++ предпочитает специализированные casts
\texttt{static\_cast}, \texttt{dynamic\_cast}, \texttt{const\_cast} и
\texttt{reinterpret\_cast} вместо C-style cast.

\begin{importantbox}[Сквозная идея]
Move и C++ casts не являются «магическими действиями».
Они меняют то, \emph{как expression участвует в type system}.
После этого обычные constructors, overload resolution и conversion rules
определяют реальную операцию.
\end{importantbox}

\subsection{Почему одного copy недостаточно}

Представим object, который владеет большим buffer или
\texttt{std::unique\_ptr<Resource>}. Copy должен создать независимое состояние:
для большого buffer это может означать allocation и копирование всех элементов,
а unique ownership вообще не допускает двух owners одного resource.

Но часто старый object после передачи больше не нужен:

\begin{cppcode}
Packet packet = make_packet();
send(/* packet больше не нужен после этого вызова */);
\end{cppcode}

Вместо глубокого copy можно \term{передать} внутренний resource новому object,
а старый оставить в корректном, но опустошённом состоянии. Это и есть главный
смысл move semantics.

\subsection{Expression и object --- не одно и то же}

Object существует в памяти в течение lifetime. Expression --- фрагмент
программы, который вычисляется и имеет type и value category.

Один и тот же object может появляться в expressions разных категорий:

\begin{cppcode}
Packet packet{"report"};

inspect(packet);            // expression packet: lvalue
inspect(std::move(packet)); // expression std::move(packet): xvalue
\end{cppcode}

Сам object \texttt{packet} не превращается в другой object. Меняется категория
expression, через которое мы к нему обращаемся.

\subsection{Карта value categories в C++17}

Удобно сначала увидеть полную схему:

\begin{center}
\begin{tabularx}{\textwidth}{L{0.22\textwidth}YY}
\toprule
Категория & Интуиция & Обычный пример \\
\midrule
lvalue & glvalue с устойчивой идентичностью & named variable \\
xvalue & glvalue, чей resource можно переиспользовать & \texttt{std::move(x)} \\
prvalue & вычисляет значение для инициализации/result & \texttt{Packet\{"x"\}} \\
glvalue & lvalue или xvalue; expression идентифицирует entity & \texttt{x}, \texttt{std::move(x)} \\
rvalue & prvalue или xvalue & temporary, \texttt{std::move(x)} \\
\bottomrule
\end{tabularx}
\end{center}

Эта схема точнее старой фразы «lvalue стоит слева от присваивания».
Например, const lvalue может вообще не быть допустимой левой частью assignment.

\begin{standardbox}[Минимальная точная модель]
В C++17 lvalue и xvalue относятся к glvalue; prvalue и xvalue относятся к
rvalue. Xvalue одновременно glvalue и rvalue.
\end{standardbox}

\subsection{Lvalue}

Named variable expression обычно lvalue:

\begin{cppcode}
std::string name{"block-12"};
std::string& ref = name;
\end{cppcode}

Lvalue удобно понимать как expression, который идентифицирует уже существующий
object/function и обычно позволяет продолжать обращаться к той же сущности.

Но наличие адреса --- плохое универсальное определение lvalue: язык задаёт
категории семантически, а не тестом «можно ли вызвать operator \&».

\subsection{Prvalue}

Temporary/value-producing expression часто prvalue:

\begin{cppcode}
std::string{"temporary"}
7 + 5
make_value()
\end{cppcode}

В C++17 prvalue тесно связан с инициализацией результата. Не нужно представлять,
что compiler обязан сначала создать отдельный временный object, а затем
обязательно скопировать его в destination: C++17 во многих случаях прямо
строит result object на нужном месте.

\subsection{Xvalue}

\term{Xvalue} --- glvalue, обозначающий object, resource которого разрешено
переиспользовать. Самый знакомый источник xvalue:

\begin{cppcode}
std::move(packet)
\end{cppcode}

Именно к xvalue обычно привязывается rvalue reference, позволяя выбрать move
overload.

\subsection{Rvalue references}

Синтаксис:

\begin{cppcode}
Packet&& ref = Packet{"temporary"};
\end{cppcode}

\texttt{Packet\&\&} --- rvalue reference. Она может привязаться к подходящему
rvalue. Такой parameter позволяет overload-у сказать:
«caller предоставляет object, состояние которого можно забрать».

Обычная пара:

\begin{cppcode}
void process(const Packet& packet); // borrow/read
void process(Packet&& packet);      // можно переиспользовать resource
\end{cppcode}

Overload resolution выбирает подходящую функцию из type и value category
argument.

\subsection{Named rvalue reference --- lvalue-expression}

Это один из самых важных моментов:

\begin{cppcode}
void relay(Packet&& packet) {
    process(packet);            // packet --- named expression => lvalue
    process(std::move(packet)); // xvalue => rvalue overload
}
\end{cppcode}

У variable \texttt{packet} declared type содержит \texttt{\&\&}, но expression
с её именем --- lvalue. Поэтому move не «распространяется автоматически»
через named variables.

\subsection{Что на самом деле делает \texttt{std::move}}

Несмотря на имя, \texttt{std::move} сам по себе не переносит bytes и не
обнуляет object. Концептуально он выполняет cast к подходящей rvalue-reference
категории:

\begin{cppcode}
// Упрощённая идея, не реализация для копирования:
static_cast<T&&>(object)
\end{cppcode}

Реальная library template учитывает reference removal. Результат
\texttt{std::move(x)} --- xvalue. После этого overload resolution может выбрать
move constructor или move assignment operator.

\begin{whybox}[Почему после \texttt{std::move} иногда всё равно происходит copy?]
\texttt{std::move} только меняет category expression. Если подходящей move
operation нет, но существует copy operation, способная принять такой argument,
compiler может выбрать copy. Поэтому фраза «std::move перемещает» технически
неточна.
\end{whybox}

\subsection{Move constructor}

Обычная форма:

\begin{cppcode}
class Packet {
public:
    Packet(Packet&& other) noexcept
        : data_{std::move(other.data_)} {}
};
\end{cppcode}

Новый Packet создаётся из rvalue другого Packet. Для unique-owned resource
move constructor обычно передаёт ownership member-у нового object.

Если member --- \texttt{unique\_ptr}, его собственная move operation делает
source smart pointer пустым.

\subsection{Move assignment}

Move assignment работает с \emph{уже существующим} destination:

\begin{cppcode}
Packet& operator=(Packet&& other) noexcept {
    if (this != &other) {
        data_ = std::move(other.data_);
    }
    return *this;
}
\end{cppcode}

Как и в copy semantics из блока 06:

\begin{cppcode}
Packet b = std::move(a); // construction
b = std::move(c);        // assignment
\end{cppcode}

--- это разные операции.

Self-move нужно учитывать в contract конкретного type. Не следует строить
общую программу на предположении, что после \texttt{x = std::move(x)} состояние
обязательно сохраняется прежним.

\subsection{Moved-from object}

После move source object продолжает существовать, пока не закончился его
lifetime. Но его состояние определяется contract типа.

Для многих standard-library types после move гарантируется состояние
\emph{valid but unspecified}: invariants сохранены, destructor и операции,
не требующие конкретного старого значения, допустимы, но прежнее содержимое
предполагать нельзя.

\begin{warningbox}[Не универсальное правило «после move всегда пусто»]
\texttt{unique\_ptr} после move действительно становится empty по своему
contract. Для произвольного \texttt{std::string} нельзя переносимо строить
логику на том, что \texttt{empty()} обязательно true после move.
\end{warningbox}

Для собственного \texttt{Packet} мы вправе определить более сильный contract:
например, moved-from packet становится пустым и может быть заново присвоен.

\subsection{Rule of Five}

В блоке 06 был Rule of Three. После появления move operations design-guideline
расширяется до \term{Rule of Five}: если class вручную определяет resource
semantics и вынужден писать один из
destructor; copy constructor; copy assignment; move constructor; move assignment.
нужно осознанно рассмотреть все пять.

Это не требование синтаксиса C++.

\subsection{Rule of Zero по-прежнему лучше}

Современный class часто не должен писать ни одну из пяти operations вручную.
Если members сами корректно управляют lifetime, compiler-generated special
members могут быть правильным решением.

Например, class с \texttt{std::string} и \texttt{std::vector} обычно получает
правильную copy/move semantics из своих members.

Class с \texttt{unique\_ptr} автоматически не становится copyable, но при
подходящих условиях может получить move operations. Не добавляйте ручной move
constructor только ради демонстрации синтаксиса в production code.

\subsection{Почему \texttt{noexcept} важен для move}

Move operation resource-owning type часто может просто передать handles и не
требует операций, способных бросить exception. Тогда её можно объявить
\texttt{noexcept}.

Standard containers при перераспределении storage должны соблюдать свои
exception-safety contracts. Для copyable type они могут предпочесть copy вместо
potentially-throwing move, если это помогает сохранить guarantee.

\begin{performancebox}[Не превращаем это в миф]
Нельзя утверждать «vector всегда вызывает move при reallocation».
Выбор зависит от properties element type и требований операции. Поэтому
\texttt{noexcept} move важен не только как документация, но и как информация,
которую library algorithm может использовать.
\end{performancebox}

\subsection{Copy elision и RVO}

Move semantics не означает, что каждый return object проходит через move
constructor.

В C++17:

\begin{cppcode}
Packet make_packet() {
    return Packet{"audit"};
}

Packet packet = make_packet();
\end{cppcode}

в соответствующем prvalue case result строится непосредственно как destination;
отдельный copy/move constructor для промежуточного temporary не требуется.

\term{NRVO} для named local:

\begin{cppcode}
Packet make_packet() {
    Packet result{"audit"};
    return result;
}
\end{cppcode}

--- разрешённая optimization, но не тем же самым обязательным случаем.

\begin{mistakebox}[Не писать \texttt{return std::move(local)} автоматически]
Принудительный \texttt{std::move} named local в return может помешать NRVO.
Обычный \texttt{return local;} чаще выражает намерение правильнее: language
rules уже умеют использовать copy elision/move там, где это предусмотрено.
\end{mistakebox}

\subsection{Move-only types как API contract}

\texttt{unique\_ptr} --- классический move-only type. Собственный object также
может быть move-only:

\begin{cppcode}
class Packet {
public:
    Packet(const Packet&) = delete;
    Packet& operator=(const Packet&) = delete;

    Packet(Packet&&) noexcept = default;
    Packet& operator=(Packet&&) noexcept = default;
};
\end{cppcode}

Это сообщает: два независимых copies не имеют смысла, зато ownership/state
можно передать.

В эксперименте \texttt{move\_ownership\_pipeline} этот contract проверяется
реальными compile-failure cases.

\subsection{Forwarding reference}

После templates из блока 09 можно понять ещё один случай \texttt{T\&\&}.
Если \texttt{T} --- deduced template parameter в форме:

\begin{cppcode}
template <typename T>
void relay(T&& value);
\end{cppcode}

то это \term{forwarding reference}. Она способна принять и lvalue, и rvalue.
При lvalue-вызове deduction и reference-collapsing rules сохраняют информацию,
что source был lvalue.

Не каждый \texttt{\&\&} --- forwarding reference:
\texttt{Packet\&\&} без deduction --- обычная rvalue reference, как и
\texttt{const T\&\&} в template.

\subsection{Reference collapsing}

Для forwarding достаточно запомнить таблицу:

\begin{center}
\begin{tabular}{ll}
\toprule
Комбинация & Результат \\
\midrule
\texttt{T\& \&} & \texttt{T\&} \\
\texttt{T\& \&\&} & \texttt{T\&} \\
\texttt{T\&\& \&} & \texttt{T\&} \\
\texttt{T\&\& \&\&} & \texttt{T\&\&} \\
\bottomrule
\end{tabular}
\end{center}

Практический смысл: наличие lvalue-reference на любом этапе «побеждает»,
поэтому forwarding template может сохранить исходную category.

\subsection{\texttt{std::forward}}

Named parameter \texttt{value} внутри relay всё равно
lvalue-expression. Простая передача потеряет rvalue category:

\begin{cppcode}
template <typename T>
void bad_relay(T&& value) {
    process(value); // всегда named expression => lvalue
}
\end{cppcode}

\texttt{std::forward<T>} восстанавливает category, которую template получил от
caller:

\begin{cppcode}
template <typename T>
void relay(T&& value) {
    process(std::forward<T>(value));
}
\end{cppcode}

Это основа \term{perfect forwarding}: intermediate generic layer не должен
без причины превращать rvalue в lvalue или наоборот.

\begin{warningbox}[Move и forward решают разные задачи]
\texttt{std::move(x)} безусловно выражает «дальше x можно рассматривать как
источник для move». \texttt{std::forward<T>(x)} условно сохраняет category,
deduced от caller. Не заменяйте один инструмент другим механически.
\end{warningbox}

\subsection{Move как cast: мост ко второй половине лекции}

Теперь видно, почему тема casts естественно следует за move.
\texttt{std::move} не выполняет resource transfer сам: его роль близка к
явному преобразованию value category. Реальный transfer делает выбранная
special member function.

В C++ есть и другие случаи, когда programmer явно говорит compiler:
«эту conversion я намеренно разрешаю». Для разных смыслов язык предоставляет
разные casts.

\subsection{Implicit conversion или явный cast?}

Implicit conversion хороша, когда преобразование естественно и contract
очевиден:

\begin{cppcode}
double value = 42; // int -> double
\end{cppcode}

Если conversion может терять информацию или её смысл важен для читателя,
явность полезна:

\begin{cppcode}
double value = 7.9;
int truncated = static_cast<int>(value);
\end{cppcode}

Cast не делает conversion автоматически безопасной. Он лишь явно фиксирует
намерение применить конкретные conversion rules.

\subsection{\texttt{static\_cast}}

\texttt{static\_cast<T>(expression)} выполняет набор conversions, которые
compiler может проверить на основе static types.

Обычные применения:
числовое преобразование, которое хотим отметить явно; conversion enum к integer и обратно при понимании диапазонов; некоторые conversions внутри class hierarchy, если correctness уже доказана другой логикой; получение обратно typed pointer из подходящего \texttt{void*}.

\begin{cppcode}
double ratio = 7.9;
int whole = static_cast<int>(ratio);

Severity severity = Severity::warning;
int code = static_cast<int>(severity);
\end{cppcode}

\begin{warningbox}[Явный cast не проверяет диапазон за вас]
Например, floating-to-integer conversion имеет собственные ограничения
representability. \texttt{static\_cast<int>(x)} не означает
«безопасно зажать x в диапазон int». Range policy нужно проектировать отдельно.
\end{warningbox}

\subsection{Static downcast и его предел}

Если \texttt{Derived} наследуется от \texttt{Base}, compiler может разрешить
определённый \texttt{static\_cast<Derived*>} из base pointer. Но такой cast не
проверяет dynamic type object во время выполнения.

Поэтому использовать static downcast можно только когда program logic уже
гарантирует concrete type. Для runtime-проверяемого polymorphic downcast нужен
\texttt{dynamic\_cast}.

\subsection{\texttt{dynamic\_cast}}

\texttt{dynamic\_cast} работает с polymorphic hierarchy и использует runtime
type information для проверяемого downcast/cross-cast.

Pointer form:

\begin{cppcode}
Report* base = find_report();

if (auto* lecture = dynamic_cast<LectureReport*>(base)) {
    lecture->print_lines();
}
\end{cppcode}

Если object не подходит target type, pointer result равен
\texttt{nullptr}.

Reference form:

\begin{cppcode}
try {
    LectureReport& lecture = dynamic_cast<LectureReport&>(report);
} catch (const std::bad_cast&) {
    // dynamic type не подходит
}
\end{cppcode}

При неудаче reference cast бросает \texttt{std::bad\_cast}.

\subsection{Когда base должен быть polymorphic}

Для runtime downcast source type должен удовлетворять требованиям
\texttt{dynamic\_cast}; обычный случай --- polymorphic base, имеющий хотя бы
одну virtual function.

Эксперимент содержит real compile-failure, где programmer пытается выполнить
dynamic downcast из non-polymorphic base.

\subsection{Нужен ли cast для обычного upcast?}

Нет. Derived pointer/reference обычно неявно преобразуется к доступному base
pointer/reference:

\begin{cppcode}
LectureReport lecture;
Report& report = lecture; // нормальный upcast
\end{cppcode}

Добавлять \texttt{dynamic\_cast<Report\&>(lecture)} только ради «явности»
не нужно: cast должен решать проблему, а не украшать код.

\subsection{\texttt{const\_cast}}

\texttt{const\_cast} меняет cv-qualification там, где language rules это
разрешают. Главный учебный случай --- различить:
1) view const, но underlying object на самом деле mutable; 2) underlying object изначально создан как const.

Первый случай:

\begin{cppcode}
int counter = 5;
const int* view = &counter;

int* mutable_view = const_cast<int*>(view);
*mutable_view = 7; // underlying object counter не const
\end{cppcode}

Второй:

\begin{cppcode}
const int limit = 10;
int* p = const_cast<int*>(&limit);
// *p = 11; // Undefined behavior: настоящий object const
\end{cppcode}

\begin{ubbox}[Cast не отменяет constness самого object]
\texttt{const\_cast} может дать type без const, но если original object
действительно const, попытка изменить его через такой access приводит к
undefined behavior.
\end{ubbox}

Практически \texttt{const\_cast} часто появляется при interop со старым API.
В нормальном собственном interface лучше сразу правильно описать const contract.

\subsection{\texttt{reinterpret\_cast}}

\texttt{reinterpret\_cast} предназначен для низкоуровневых conversions, где
programmer явно работает с representation/pointer mapping.

Пример эксперимента делает только контролируемый pointer--integer round trip:

\begin{cppcode}
int value = 42;
int* pointer = &value;

std::uintptr_t bits =
    reinterpret_cast<std::uintptr_t>(pointer);
int* restored =
    reinterpret_cast<int*>(bits);
\end{cppcode}

На implementation, предоставляющей \texttt{std::uintptr\_t}, этот integer type
способен хранить преобразованное object pointer representation. Mapping в
integer implementation-defined; превращать число в переносимый «адрес в
формате стандарта» нельзя.

\begin{warningbox}[reinterpret не создаёт object другого типа]
Получить pointer другого type и сразу разыменовать его --- совсем другой вопрос:
вступают в силу lifetime и aliasing rules. \texttt{reinterpret\_cast} не делает
такой доступ автоматически корректным.
\end{warningbox}

\subsection{C-style cast скрывает намерение}

Запись:

\begin{cppcode}
Target value = (Target)source;
\end{cppcode}

может использовать несколько разных C++ conversion mechanisms. Читателю
сложнее понять, что именно programmer хотел разрешить, а поиск опасных
low-level conversions становится труднее.

Например, C-style pointer cast способен одновременно убрать const qualification
в коде, который не прошёл бы как обычная implicit conversion.

Поэтому production-style курса:
сначала попытаться обойтись без cast; если cast действительно нужен --- выбрать самый узкий named C++ cast; C-style casts оставлять для чтения legacy/C-oriented code, а не как основной стиль C++.

\subsection{Четыре casts рядом}

\begin{center}
\begin{tabularx}{\textwidth}{L{0.22\textwidth}YY}
\toprule
Cast & Главный смысл & Проверка/риск \\
\midrule
\texttt{static\_cast} & compile-time known conversion & не доказывает runtime dynamic type \\
\texttt{dynamic\_cast} & checked polymorphic cast & runtime check; pointer null/reference exception \\
\texttt{const\_cast} & изменить cv-qualification & mutation настоящего const object может быть UB \\
\texttt{reinterpret\_cast} & low-level reinterpretation/mapping & lifetime, aliasing, portability \\
\bottomrule
\end{tabularx}
\end{center}

\subsection{Cast не должен маскировать design problem}

Если code постоянно делает \texttt{dynamic\_cast} по длинной цепочке concrete
types, возможно, common virtual interface спроектирован слишком слабо.

Если code постоянно использует \texttt{const\_cast}, возможно, constness API
описана неверно.

Если application logic построена на \texttt{reinterpret\_cast}, нужно проверить,
не скрывается ли за этим неподходящее representation или нарушение ownership.

Cast --- точечный инструмент, не замена хорошему type design.

\subsection{Эксперимент 1: move\_ownership\_pipeline}

Каталог:
\filepath{examples/12\_move\_semantics\_and\_casts/01\_move\_ownership\_pipeline}.

Move-only \texttt{Packet} владеет payload через \texttt{unique\_ptr}. Lab
проверяет:
1) copy construction и copy assignment запрещены; 2) move construction передаёт resource и оставляет source empty по contract именно этого Packet; 3) move assignment делает то же для уже существующего destination; 4) named variable выбирает lvalue overload, temporary и \texttt{std::move} --- rvalue overload; 5) named rvalue-reference variable снова lvalue-expression; 6) forwarding function с \texttt{std::forward} сохраняет category; 7) prvalue return в C++17 не требует промежуточного move в проверяемом case.

\begin{shellcode}
cd examples/12_move_semantics_and_casts/01_move_ownership_pipeline
./run_checks.sh
\end{shellcode}

Два negative compile cases подтверждают отдельно:
copy move-only Packet запрещён, а const \texttt{unique\_ptr} нельзя передать
обычным move construction, потому что его resource нельзя забрать из const
owner.

\subsection{Эксперимент 2: casts\_lab}

Каталог:
\filepath{examples/12\_move\_semantics\_and\_casts/02\_casts\_lab}.

Positive executable проверяет:
явную numeric conversion через \texttt{static\_cast}; успешный и неуспешный pointer \texttt{dynamic\_cast}; exception от reference \texttt{dynamic\_cast}; безопасную mutation через \texttt{const\_cast}, когда original object действительно non-const; pointer--integer round trip через \texttt{reinterpret\_cast} на текущей implementation.

Два negative compile cases проверяют const access и попытку runtime downcast из
non-polymorphic base. Ещё один compile-only case показывает, что C-style cast
способен убрать const в source code --- программа его только компилирует и
\emph{не выполняет опасную запись}.

\begin{shellcode}
cd examples/12_move_semantics_and_casts/02_casts_lab
./run_checks.sh
\end{shellcode}

\subsection{Стоимость}

Move часто превращает глубокое копирование resource в перенос нескольких
handles/pointers, но это свойство конкретного type. Для small trivially copied
value move может не давать выигрыша вообще.

\texttt{static\_cast} и \texttt{const\_cast} как языковые conversions сами по
себе не означают обязательный runtime dispatch.
\texttt{dynamic\_cast} для down/cross cast требует runtime type check и может
иметь ненулевую стоимость.
\texttt{reinterpret\_cast} тоже не «дорогая функция», но его
семантическая цена --- portability и риск нарушить низкоуровневые правила.

\begin{performancebox}[Сначала корректная semantics]
Главный вопрос блока не «что быстрее: copy или move?» и не «какой cast быстрее?».
Сначала нужен корректный ownership/type contract. Измерение имеет смысл только
после того, как обе сравниваемые реализации решают одну задачу корректно.
Поэтому отдельный benchmark в блоке не нужен.
\end{performancebox}

\subsection{Итоговая схема API}

Передача аргумента:
\texttt{const T\&} --- borrow read-only object; \texttt{T\&} --- borrow mutable object; \texttt{T} by value --- функция получает своё value; caller может copy или move в него; \texttt{T\&\&} --- конкретный API принимает rvalue и может переиспользовать состояние; \texttt{T\&\&} с deduced template parameter --- forwarding reference, если выполнены соответствующие условия.

Явная conversion:
сначала проверить, нельзя ли написать interface без cast; \texttt{static\_cast} --- обычная compile-time conversion; \texttt{dynamic\_cast} --- runtime-checked polymorphic conversion; \texttt{const\_cast} --- узкая работа с cv-qualification; \texttt{reinterpret\_cast} --- низкоуровневая boundary.

\begin{importantbox}[Главный вывод]
\texttt{std::move} не перемещает object: он разрешает рассматривать expression
как источник для move. Реальное изменение состояния выполняет выбранная
operation типа. Точно так же named C++ casts не «делают всё безопасным», а
выражают конкретное намерение conversion и позволяют отдельно рассуждать о его
preconditions.
\end{importantbox}

В следующем блоке мы разберём type deduction, \texttt{auto}, \texttt{decltype},
structured bindings, CTAD, lambdas, \texttt{std::function} и
\texttt{std::optional}. Они используют уже построенную здесь модель expressions,
types и value categories.
